% ..............................................................................
% Slide 6 -- Method: self-supervised learning with JEPA
% Source content: contents/slide-06-method.md  (bullets are FINAL -- render as-is)
% ------------------------------------------------------------------------------
\section{Method}

\begin{frame}[t]{Self-supervised learning with JEPA}
% ---- Diagram: full width across the top ----------------------------------
\centering
\includegraphics[width=\linewidth,height=0.5\textheight,keepaspectratio]%
{contents/images/jepa.png}

\vfill
% ---- Bullets: lower third, two compact columns ---------------------------
{\small
\begin{columns}[t]
\begin{column}{0.5\textwidth}
\begin{itemize}\itemsep=0.25em
\item Self-supervision: learn from raw, \textbf{unlabeled} data by predicting
      hidden parts of it
\item The engine behind modern AI: LLMs learn by predicting the next word
\item \textbf{JEPA}: predict in \textbf{latent space, not pixels} $\to$ abstract
      structure, ignores pixel noise
      \cite{lecun2022path, assran2023ijepa, bardes2024vjepa}
\end{itemize}
\end{column}
\begin{column}{0.5\textwidth}
\begin{itemize}\itemsep=0.25em
\item My model: context encoder + EMA target encoder \cite{grill2020byol} +
      predictor, with variance-covariance regularization \cite{bardes2022vicreg}
      $\to$ pretrain on unlabeled video, then a light segmentation head
\item \textbf{My contribution}: adapt JEPA-style self-supervision to
      calcium-imaging segmentation
\end{itemize}
\end{column}
\end{columns}
}
\end{frame}
